% Solutions of 12-program-correctness.
\section*{Chapter \ref{ch:b1:program-correctness} --- Program Correctness}

\begin{solution}{exo:b1:program-correctness:1}
In Python, for $b > 0$ the remainder \texttt{a \% b} lies in $[0, b)$, so $a - (a \bmod b)$ is the largest multiple of
$b$ that is at most $a$. (a) Met. (b) Met, since $0 \leq a \bmod b < b$. (c) Not met: for $b < 0$ the remainder lies in
$(b, 0]$, and \code{f(5, -3)} returns 6, above $a$. (d) Met, for negative $a$ too: Python's remainder takes the sign of
$b$.
\end{solution}

\begin{solution}{exo:b1:program-correctness:2}
At the five tests, $(i, m)$ is $(1, 3)$, $(2, 9)$, $(3, 9)$, $(4, 9)$ and $(5, 9)$, and the loop exits at the last. Each
time $m$ is the maximum of $xs[0{:}i]$: of $[3]$, $[3, 9]$, $[3, 9, 2]$, $[3, 9, 2, 9]$ and the whole list.
\end{solution}

\begin{solution}{exo:b1:program-correctness:3}
(a) Stops for even $n$, with variant $n$; for odd $n$ it passes 1, then $-1, -3, \dots$ and never stops. (b) Stops:
$n$ is a variant, halved while $n > 1$. (c) Stops only for $n = 0$ (no iteration) and $n = 1$; from $n = 2$ the values
are $3, 2, 3, 2, \dots$, and from any $n \geq 2$ the loop alternates between an even and an odd number for ever.
(d) Stops: since 3 and 7 are coprime, $n, n + 3, \dots, n + 18$ meet every remainder modulo 7; the variant is the
least $k \geq 0$ with $7 \mid n + 3k$, at most 6, which drops by one at each iteration.
\end{solution}

\begin{solution}{exo:b1:program-correctness:4}
\omoutput[2-3]{code/bachelor-1/12-program-correctness/outputs/exo_dups.txt}{The leftmost positions.}
In $[1, 1, 1]$: the first 1 is at 0; 2 would go at the end (3); 0 at the start (0). In $[1, 2, 2, 2, 3]$: the first 2
is at 1, the 3 at 4; in $[4]$, 4 is at 0.
\end{solution}

\begin{solution}{exo:b1:program-correctness:5}
On a closed interval the variant is the number of candidates, $hi - lo + 1$. When $lo < hi$, $mid < hi$ and
\code{hi = mid} decreases it; but when $lo = hi = mid$ and $a[mid] > x$, \code{hi = mid} changes nothing, and the loop
repeats the same state. This happens whenever $x$ is absent and the search ends next to a larger element: in
$[2, 5, 8]$, for $x = 1$ ($lo = hi = 0$), and also for $x = 3$ or $x = 6$. The fix is \code{hi = mid - 1}, with the
invariant ``every element of $a[{:}lo]$ is $< x$ and every element of $a[hi + 1{:}]$ is $> x$'': $a[mid] > x$ and the
list is sorted, so $a[mid{:}]$ is $> x$; and the variant $hi - lo + 1$ now decreases in all three branches. The
half-open version of \cref{lst:b1:program-correctness:bsearch} avoids mixing the two conventions.
\end{solution}

\begin{solution}{exo:b1:program-correctness:6}
Preservation fails. With \code{a[k] = a[j]} the element at $k$, which belongs to the region $\geq p$, is overwritten and
lost, and the small $a[j]$ stays where it was, inside the region $a[k{:}j]$, which is no longer $\geq p$. The result is
not a permutation (7 and 9 are lost, 1 and 4 appear twice), and it is not even partitioned: 1 and 4 come after the
pivot 5. The test of \cref{prop:b1:program-correctness:partition}'s two conclusions catches both at once.
\end{solution}

\begin{solution}{exo:b1:program-correctness:7}
\omcode{code/bachelor-1/12-program-correctness/exercises/rightmost.py}{4}{14}{One past the last occurrence}
The proof is that of \cref{prop:b1:program-correctness:bsearch}, with $\leq$ in place of $<$: if $a[mid] \leq x$, all
of $a[{:}mid + 1]$ is $\leq x$; otherwise all of $a[mid{:}]$ is $> x$. The occurrences of $x$ are exactly
$a[\mathit{leftmost}{:}\mathit{rightmost}]$, so their number is the difference, as the output of
\cref{lst:b1:program-correctness:bsearch} showed for $x = 3$. The test compares both functions with \code{bisect_left}
and \code{bisect_right} on random sorted lists with many duplicates.
\end{solution}

\begin{solution}{exo:b1:program-correctness:8}
Let $n = \text{len}(xs) \geq 1$. Invariant: $1 \leq i \leq n$, $m \in xs[0{:}i]$ and $m \geq x$ for every $x$ of
$xs[0{:}i]$. Initialisation: $i = 1$, $m = xs[0]$, the maximum of $[xs[0]]$. Preservation: with $i < n$, if $xs[i] > m$
the new $m$ is $xs[i]$, larger than every element of $xs[0{:}i]$; otherwise $m \geq xs[i]$ already; either way $m$ is
the maximum of $xs[0{:}i + 1]$, and $i + 1 \leq n$. Conclusion: at the exit $i \geq n$ and $i \leq n$, so $i = n$ and
$m$ is the maximum of the whole list. Termination: $n - i$ is a natural number that decreases by one.
\end{solution}

\begin{solution}{exo:b1:program-correctness:9}
Let $L = hi - lo \geq 1$ and $mid = lo + \lfloor L/2 \rfloor$. After \code{hi = mid} the length is $\lfloor L/2 \rfloor$;
after \code{lo = mid + 1} it is $L - \lfloor L/2 \rfloor - 1 = \lceil L/2 \rceil - 1 \leq \lfloor L/2 \rfloor$. So after
$k$ iterations the length is at most $\lfloor n/2^k \rfloor$, which is 0 for $k = \lfloor \log_2 n \rfloor + 1$, since
then $2^k > n$: the loop has stopped. The bound is reached when $x$ is below every element: every step takes
\code{hi = mid}, and the length goes $n, \lfloor n/2 \rfloor, \lfloor n/4 \rfloor, \dots$, which stays positive for
$\lfloor \log_2 n \rfloor + 1$ values. The test checks the worst case for every $n < 300$.
\end{solution}

\begin{solution}{exo:b1:program-correctness:10}
On $[2, 2, 0]$ the flag makes 3 swaps (two of them swap an element with itself), where one swap of the first and
last elements sorts the list; on the list of \cref{out:b1:program-correctness:flag}, 6 swaps where 3 suffice.
\omcode{code/bachelor-1/12-program-correctness/exercises/min_swaps.py}{4}{15}{The fewest swaps}
Count $c[i][j]$, the elements of value $i$ in the region where $j$ belongs. A swap moves at most two elements to their
place, and does so only for a pair $i \leftrightarrow j$ with $c[i][j] > 0$ and $c[j][i] > 0$; take as many such
swaps as possible. The misplaced elements left form cycles $0 \to 1 \to 2 \to 0$ (or the reverse), in equal numbers
$t$ in each of the three counts, and a cycle of three needs two swaps, each placing at most one element beyond the
pairs. The minimum is $\text{pairs} + 2t$, with $3t$ the elements left. The test compares it with a breadth-first search
over all orders reached by swaps, for every list of 0, 1 and 2 of length at most 6, and checks that the flag never does
better.
\end{solution}

\begin{solution}{exo:b1:program-correctness:11}
\omoutput[2-2]{code/bachelor-1/12-program-correctness/outputs/exo_mystery.txt}{Three values.}
It computes $x^n$: $3^{13} = 1\,594\,323$, $2^{10} = 1024$, $7^0 = 1$. \omterm{def:b1:program-correctness:specification}{Specification}: pre $n \geq 0$, post the result
is $x^n$. Invariant: $\mathit{acc} \cdot \mathit{base}^e = x^n$ and $e \geq 0$. Initially $1 \cdot x^n$. If $e$ is
odd, $\mathit{acc} \cdot \mathit{base} \cdot (\mathit{base}^2)^{(e - 1)/2} = \mathit{acc} \cdot \mathit{base}^e$; if
$e$ is even, $\mathit{acc} \cdot (\mathit{base}^2)^{e/2} = \mathit{acc} \cdot \mathit{base}^e$. At the exit $e = 0$ and
$\mathit{acc} = x^n$. Variant: $e$, halved at each step. The number of iterations is the number of binary digits of $n$,
$\lfloor \log_2 n \rfloor + 1$ for $n \geq 1$, and 0 for $n = 0$: 4 for $n = 13$. The test compares with \code{x ** n}.
\end{solution}

\begin{solution}{exo:b1:program-correctness:12}
The first proof is right, for $n \geq 0$: $s + i + 1 = (i + 1)(i + 2)/2$ preserves the first part, $i < n$ before the
step gives $i + 1 \leq n$ after it, and the exit gives $i = n$. The second is a correct proof of \emph{partial}
correctness only: at the exit $i = n$ by the loop condition alone, and the invariant ``$i$ is even'' plays no part. What
is missing is termination, and it is false for odd $n$ (and negative $n$), where $i$ jumps over $n$ for ever. The
invariant should be ``$i$ is even and $i \leq n$'', which needs the \omterm{def:b1:program-correctness:specification}{precondition} ``$n$ even and $n \geq 0$''; then
$n - i$ is a variant.
\end{solution}

\begin{solution}{pb:b1:program-correctness:1}
\textbf{1.} Existence: $\{r \in \mathbb{N} : r^2 \leq n\}$ contains 0 and is bounded by $n + 1$; its largest element
$r$ has $(r + 1)^2 > n$. Uniqueness: if $r < r'$ both satisfied it, then $r + 1 \leq r'$ and $(r + 1)^2 \leq r'^2 \leq
n$, a contradiction.

\textbf{2.} Compute $r \cdot r$ and $(r + 1) \cdot (r + 1)$ and compare them with $n$.

\textbf{3.} A \code{float} has a 53-bit significand: above $2^{53} \approx 9 \times 10^{15}$ not every integer is
representable, and $n$ is rounded before the root is taken. For $n = 2^{64} - 1$, \code{math.floor(math.sqrt(n))}
gives 4\,294\,967\,296, one too many.

\textbf{4.} \omterm{def:b1:program-correctness:specification}{Precondition}: $n$ is an integer, $n \geq 0$. \omterm{def:b1:program-correctness:specification}{Postcondition}: the result is $(r, k)$ with $r^2 \leq n < (r +
1)^2$ and $k$ the number of iterations made.

\textbf{5.} \code{assert r * r <= n < (r + 1) * (r + 1)}, which needs neither program nor library.

\textbf{6.} To compare the methods in Part IV; the test bounds them as well.

\textbf{7.} $0^2 = 0 \leq n$ and $n < (n + 1)^2$.

\textbf{8.} In an iteration $hi \geq lo + 2$, so $lo + 1 \leq mid \leq hi - 1$. If $mid^2 \leq n$, \code{lo = mid}
keeps $lo^2 \leq n$; otherwise $n < mid^2$ and \code{hi = mid} keeps $n < hi^2$.

\textbf{9.} $hi - lo$ is a natural number ($lo^2 \leq n < hi^2$ gives $lo < hi$) and decreases strictly, since $mid$
lies strictly between $lo$ and $hi$; the loop runs only while it exceeds 1.

\textbf{10.} At the exit $lo < hi \leq lo + 1$, so $hi = lo + 1$, and the invariant reads $lo^2 \leq n < (lo + 1)^2$.

\textbf{11.} A step leaves $mid - lo = \lfloor L/2 \rfloor$ or $hi - mid = \lceil L/2 \rceil$, so after $k$ steps the
length is at most $\lceil (n + 1)/2^k \rceil$, which is 1 once $2^k \geq n + 1$: at most $\lceil \log_2(n + 1) \rceil$
steps. The table reaches the bound except for 16, 17, 99 and $2^{32}$, one below.

\textbf{12.} $n + 1$ is 0 for $n = 2^{64} - 1$; $lo + hi$ can exceed $2^{64}$; and $mid \cdot mid$ overflows once $mid
\geq 2^{32}$. Start with $hi = 2^{32}$ (the root of a 64-bit number is below it, so $n < hi^2$ holds), compute $mid = lo
+ (hi - lo)/2$, and test $mid \leq n / mid$ (for $mid \geq 1$) instead of $mid^2 \leq n$.

\textbf{13.} Write $n/r = q + f$ with $q = \lfloor n/r \rfloor$ and $0 \leq f < 1$, and $m = r + q$. If $m$ is even,
$\lfloor (m + f)/2 \rfloor = m/2$; if $m$ is odd, $(m + f)/2 = (m - 1)/2 + (1 + f)/2$ with $(1 + f)/2 < 1$, so the floor is
$(m - 1)/2$. Either way it is $\lfloor m/2 \rfloor$.

\textbf{14.} $(r + n/r)/2 - \sqrt{n} = (\sqrt{r} - \sqrt{n/r})^2/2 \geq 0$. Taking floors, the new $r$ is at least
$\lfloor \sqrt{n} \rfloor = s$. Initially $r = n \geq s$, since $s \leq s^2 \leq n$ for $s \geq 1$.

\textbf{15.} $r^2 > n$ gives $n/r < r$, so $(r + n/r)/2 < r$, and its floor is below $r$.

\textbf{16.} $r$ is a natural number that decreases strictly at each iteration: the loop stops. At the exit $r^2 \leq
n$, so $r \leq s$, the largest number whose square is at most $n$; with the invariant $r \geq s$, $r = s$.

\textbf{17.} The first step computes \code{n // r} with $r = n = 0$: a division by zero.

\textbf{18.} The values can oscillate between $s$ and $s + 1$: for $n = 15$, $r = 4$ gives 3, and $r = 3$ gives 4, and
so on. The test $r^2 > n$ stops at 3.

\textbf{19.} The bound is reached for 0, 1, 2, 15, $10^6$, $2^{32} - 1$, $10^{18}$ and $2^{64} - 1$; bisection stays
one below for 16, 17, 99 and $2^{32}$, where the intervals happen to shrink a little faster than half.

\textbf{20.} While $r$ is far above $\sqrt{n}$, $n/r$ is small beside $r$, and the new $r$ is about $r/2$: from $2^{64}$
to near $2^{32}$ takes about 32 steps. Near $\sqrt{n}$ the error is roughly squared at each step (it becomes about
$e^2/(2r)$), so a handful of steps finish: $32 + 5 = 37$.

\textbf{21.} Start at $2^{\lceil b/2 \rceil}$ with $b$ the number of binary digits of $n$: since $n < 2^b$, $\sqrt{n} <
2^{b/2}$, so it is at least $s$. Only the initialisation of the invariant must be proved again; preservation,
conclusion and termination do not depend on the start.

\textbf{22.} Random 64-bit values almost never land on perfect squares, their neighbours or tiny values, where the
boundary errors live; the exhaustive check below 2000 covers them all.

\textbf{23.} For $n = 16$: $r$ goes $16, 8, 5, 4$, then $4 \cdot 4 \geq 16$ still holds and $r$ stays 4 for ever. The
variant fails when $r^2 = n$: the new $r$ equals $r$. The test would never finish: the
exhaustive loop hangs at $n = 1$, where $r = 1$ stays 1.

\textbf{24.} \emph{Newton's integer square root is \omterm{def:b1:program-correctness:correctness}{totally correct}}: for every $n \geq 0$, \code{isqrt_newton(n)}
terminates and returns the $r$ with $r^2 \leq n < (r + 1)^2$ (questions 13--16). On the table it needs 0, 0, 1, 3, 3,
3, 5, 12, 20, 19, 34 and 37 iterations, against bisection's 0, 1, 2, 4, 4, 4, 6, 20, 32, 32, 60 and 64, which never
exceed $\lceil \log_2(n + 1) \rceil$.
\end{solution}
